\chapter{2025年上海卷（秋）}

\section{填空题}

\begin{problem}
已知全集 \(U = \{x \mid 2 \le x \le 5, x \in \R\}\)，集合 \(A = \{x \mid 2 \le x < 4, x \in \R\}\)，则 \(\bar{A} =\)\fillinblank{}.
\end{problem}

\begin{answer}
\(\bar{A} = \{x \mid 4 \le x \le 5, x \in \R\}=[4,5]\)
\end{answer}

\begin{solution}
根据补集定义，\(\bar{A} = \{x \mid x \in U, x \notin A\}\). 又 \(U=\{x\mid2\le x\le5\}\)，所以 \(\bar{A}=\{x\mid4\le x\le5\}\).
\end{solution}

\begin{problem}
设 \(x\in\R\)，不等式 \(\frac{x-1}{x-3}<0\) 的解集为\fillinblank{}.
\end{problem}

\begin{answer}
\((1,3)\)
\end{answer}

\begin{solution}
由 \(\frac{x-1}{x-3}<0\)，得 \(1<x<3\).
\end{solution}

\begin{problem}
已知等差数列 \(\{a_n\}\) 的首项 \(a_1 = -3\)，公差 \(d = 2\)，则该数列的前 \(6\) 项和为\fillinblank{}.
\end{problem}

\begin{answer}
\(12\)
\end{answer}

\begin{solution}
由等差数列求和公式，\(S_6 = \frac{6(a_1+a_6)}{2}=3(a_1+a_6)\)，其中 \(a_6=a_1+5d=7\)，故 \(S_6=3(-3+7)=12\).
\end{solution}

\begin{problem}
在二项式 \((2x-1)^5\) 的展开式中，\(x^3\) 的系数为\fillinblank{}.
\end{problem}

\begin{answer}
\(80\)
\end{answer}

\begin{solution}
二项式 \((2x-1)^5\) 中 \(x^3\) 的系数对应 \(r=2\)，\(\mathrm{T}_{3}=\mathrm C_5^2(2)^3(-1)^2=80\).
\end{solution}

\begin{problem}
函数 \(y=\cos x\) 在 \([-\pi/2,\pi/4]\) 上的值域为\fillinblank{}.
\end{problem}

\begin{answer}
\([0,1]\)
\end{answer}

\begin{solution}
函数 \(y=\cos x\) 在 \([-\pi/2,0]\) 上递增，在 \([0,\pi/4]\) 上递减. 并且 \(y(-\pi/2)=0\)，\(y(0)=1\)，\(y(\pi/4)=\sqrt2/2\). 故值域为 \([0,1]\).
\end{solution}

\begin{problem}
已知随机变量 \(X\) 的分布为
\(\begin{pmatrix}5&6&7\\0.2&0.3&0.5\end{pmatrix}\)，则期望 \(E(X)=\)\fillinblank{}.
\end{problem}

\begin{answer}
\(6.3\)
\end{answer}

\begin{solution}
由题意，\(E(X)=5\times0.2+6\times0.3+7\times0.5=1+1.8+3.5=6.3\).
\end{solution}

\begin{problem}
如图，在正四棱柱 \(ABCD-A_1B_1C_1D_1\) 中，\(AC=4\sqrt2\)，\(BD_1=9\)，则该正四棱柱的体积为\fillinblank{}.

\bitmapfigure[max width=6.0cm]{img/2025/shanghai/q7_fig1.png}
\end{problem}

\begin{answer}
\(112\)
\end{answer}

\begin{solution}
因为 \(AC=4\sqrt2\)，且 \(ABCD\) 为正方形，所以 \(BA=4\)，从而 \(BD=4\sqrt2\). 又 \(BD_1=9\)，所以 \(BB_1^2+BD^2=9^2\)，得 \(BB_1=7\). 底面积为 \(BA\cdot BC=16\)，体积 \(=16\times7=112\).
\end{solution}

\begin{problem}
设 \(a\)，\(b>0\)，且 \(a>0\)，若 \(a+\frac1b=1\)，则 \(b+\frac1a\) 的最小值为\fillinblank{}.
\end{problem}

\begin{answer}
\(4\)
\end{answer}

\begin{solution}
由题设 \(a+\frac1b=1\)，得 \(0<a<1\)，且 \(b=\frac1{1-a}\). 因此
\(b+\frac1a=\frac1{1-a}+\frac1a=\frac1{a(1-a)}\ge4\).
当 \(a=\frac12\)，\(b=2\) 时取等号，故最小值为4.
\end{solution}

\begin{problem}
\(4\) 个家长和 \(2\) 个儿童去郊游爬山，\(6\) 个人需要排成一条队列，要求队列的头和尾均是家长，则不同的排列个数有\fillinblank{}.
\end{problem}

\begin{answer}
\(288\)
\end{answer}

\begin{solution}
先安排首尾两位家长，有 \(\mathrm A_4^2=12\) 种；再安排中间 \(4\) 人，有 \(\mathrm A_4^4=24\) 种，故共 \(12\times24=288\) 种.
\end{solution}

\begin{problem}
\(i\) 为虚数单位，若复数 \(z\) 满足 \(z^2=(\bar z)^2\)，且 \(|z|\le1\)，则 \(|z-2-3\i|\) 的最小值是\fillinblank{}.
\end{problem}

\begin{answer}
\(2\sqrt2\)
\end{answer}

\begin{solution}
设 \(z=a+b\i\)（\(a\)，\(b\in\R\)），则 \(\bar z=a-b\i\). 由
\(z^2=a^2+2ab\i-b^2\)，\((\bar z)^2=a^2-2ab\i-b^2\)，
可得 \(4ab\i=0\)，所以 \(ab=0\).
因此 \(a=0\) 或 \(b=0\)，即复数 \(z\) 对应的点 \(Z(a,b)\) 在线段 \([-1,1]\) 的实轴部分或线段 \([-1,1]\) 的虚轴部分上.

设 \(E(2,3)\)，则 \(|z-2-3\i|=|ZE|\).

当 \(b=0\) 时，\(|ZE|^2=(a-2)^2+9\)（\(a\in[-1,1]\)）. 因为 \(a\le1\)，所以当 \(a=1\) 时取最小值，此时 \(|ZE|_{\min}=\sqrt{(1-2)^2+9}=\sqrt{10}\).

当 \(a=0\) 时，\(|ZE|^2=4+(b-3)^2\)（\(b\in[-1,1]\)）. 因为 \(b\le1\)，所以当 \(b=1\) 时取最小值，此时
\(
|ZE|_{\min}=\sqrt{4+(1-3)^2}=2\sqrt2
\)
比较 \(\sqrt{10}\) 与 \(2\sqrt2\)，得 \(2\sqrt2<\sqrt{10}\).
故 \(|z-2-3\i|\) 的最小值为 \(2\sqrt2\).
\end{solution}

\begin{problem}
小申同学观察发现，生活中有时候影子可以完全投射在斜面上. 某斜面上有两根长为 \(1\) 米且垂直于水平面放置的杆子，与斜面的接触点分别为 \(A\)，\(B\)，它们在阳光的照射下呈现出影子，阳光可视为平行光，其中 \(A\) 处杆子的影子在水平面上，长度为 \(0.4\) 米，\(B\) 处杆子的影子完全在斜面上，长度为 \(0.45\) 米. 则斜面的坡角 \(\theta=\)\fillinblank{}（结果精确到 \(0.01^\circ\)）.

\bitmapfigure[max width=6.0cm]{img/2025/shanghai/q11_fig1.png}
\end{problem}

\begin{answer}
\(12.58^\circ\)
\end{answer}

\begin{solution}
先由 \(A\) 点杆子算出太阳光线与水平面的夹角 \(x\)，有 \(\tan x=\frac{1}{0.4}=2.5\).
在包含 \(B\) 点杆子和阳光方向的竖直截面内，设 \(C\) 为 \(B\) 点杆顶，\(E\) 为影子在斜面上的终点，过 \(E\) 作水平线交 \(B\) 点杆所在竖线于 \(D\). 设 \(BD=y\)，则 \(CD=1+y\)，且 \(\angle CED=x\)，所以 \(ED=\frac{1+y}{2.5}\).
由 \(CE=0.45\) 及勾股定理得
\(
y^2+\left(\frac{1+y}{2.5}\right)^2=0.45^2
\)
解得 \(y\approx0.098\).
故
\(
\theta=\arcsin\frac{y}{0.45}\approx12.58^\circ
\)
\end{solution}

\begin{problem}
已知函数 \(f(x)\) 满足：当 \(x>0\) 时 \(f(x)=1\)，当 \(x=0\) 时 \(f(x)=0\)，当 \(x<0\) 时 \(f(x)=-1\). \(\bs{a}\)，\(\bs{b}\)，\(\bs{c}\) 分别为平面内三个不同的单位向量. 若 \(f(\bs{a}\cdot\bs{b})+f(\bs{b}\cdot\bs{c})+f(\bs{c}\cdot\bs{a})=0\)，则 \(|\bs{a}+\bs{b}+\bs{c}|\) 的可取值范围是\fillinblank{}.
\end{problem}

\begin{answer}
\((1,\sqrt5)\)
\end{answer}

\begin{solution}
由题设，\(f(\bs{a}\cdot\bs{b})\)，\(f(\bs{b}\cdot\bs{c})\)，\(f(\bs{c}\cdot\bs{a})\in\{-1,0,1\}\). 若三者都等于 \(0\)，则
\(\bs{a}\cdot\bs{b}=0\)，\(\bs{b}\cdot\bs{c}=0\)，\(\bs{c}\cdot\bs{a}=0\).
这说明 \(\bs{a}\)，\(\bs{b}\)，\(\bs{c}\) 两两垂直. 但 \(\bs{a}\)，\(\bs{b}\)，\(\bs{c}\) 是平面内三个不同的单位向量，不可能两两垂直，矛盾.
故 \(f(\bs{a}\cdot\bs{b})\)，\(f(\bs{b}\cdot\bs{c})\)，\(f(\bs{c}\cdot\bs{a})\) 必为 \(-1\)，\(0\)，\(1\) 的一个排列. 不妨设
\(f(\bs{a}\cdot\bs{b})=1\)，\(f(\bs{b}\cdot\bs{c})=0\)，\(f(\bs{c}\cdot\bs{a})=-1\).
则
\(\bs{a}\cdot\bs{b}>0\)，\(\bs{b}\cdot\bs{c}=0\)，\(\bs{c}\cdot\bs{a}<0\).
不妨设 \(\bs{b}=(1,0)\)，\(\bs{c}=(0,1)\)，\(\bs{a}=(\cos\theta,\sin\theta)\)，则 \(\theta\in(3\pi/2,2\pi)\). 此时
\(
|\bs{a}+\bs{b}+\bs{c}|=\sqrt{(1+\cos\theta)^2+(1+\sin\theta)^2}=\sqrt{3+2\sqrt2\sin\left(\theta+\frac{\pi}{4}\right)}
\)
又有 \(\sin\left(\theta+\frac{\pi}{4}\right)\in\left(-1,\frac{\sqrt2}{2}\right)\)，故
\(
|\bs{a}+\bs{b}+\bs{c}|\in(1,\sqrt5)
\)
\end{solution}

\section{单选题}

\begin{problem}
已知事件 \(A\)，\(B\) 相互独立，若 \(P(A)=P(B)=\frac{1}{2}\)，则 \(P(A\cap B)\) 的值为（\quad）

\choices
    {\(0\)}
    {\(\frac{1}{4}\)}
    {\(\frac{1}{2}\)}
    {\(1\)}
\end{problem}

\begin{answer}
B
\end{answer}

\begin{solution}
因为 \(A\)，\(B\) 相互独立，所以 \(P(A\cap B)=P(A)P(B)=\frac{1}{2}\times\frac{1}{2}=\frac{1}{4}\). 故选 B.
\end{solution}

\begin{problem}
已知 \(a>0\) 且 \(a\ne1\)，\(s\in\R\)，下列各项中，能推出 \(a^s>a\) 的是（\quad）

\choices
    {\(a>1\) 且 \(s>0\)}
    {\(a>1\) 且 \(s<0\)}
    {\(0<a<1\) 且 \(s>0\)}
    {\(0<a<1\) 且 \(s<0\)}
\end{problem}

\begin{answer}
D
\end{answer}

\begin{solution}
因为
\(
a^s>a\iff a^{s-1}>1=a^0
\)
若 \(0<a<1\)，\(a^x\) 在 \(\R\) 上严格减函数，因此要使 \(a^{s-1}>1\) 需 \(s-1<0\). 故 \(s<1\). 又 \(0<a<1\) 时 \(a^0=1\) 为最大值，故该条件与 \(a\in(0,1)\)，\(s<0\) 对应的 D 正确.
\end{solution}

\begin{problem}
已知 \(A(0,1)\)，\(B(1,2)\)，\(C\) 在 \(\Gamma:x^2-y^2=1\ (x\ge1,y\ge0)\) 上，则 \(\triangle ABC\) 的面积（\quad）

\choices
    {有最大值，但没有最小值}
    {没有最大值，但有最小值}
    {既有最大值，也有最小值}
    {既没有最大值，也没有最小值}
\end{problem}

\begin{answer}
A
\end{answer}

\begin{solution}
设 \(C(a,b)\in\Gamma\)，则 \(a=\sqrt{b^2+1}\)，\(AB: y=x-1\). \(\triangle ABC\) 的高为 \(h=\frac{|a-b+1|}{\sqrt2}\).
设 \(f(b)=\sqrt{b^2+1}-b\)（\(b\ge0\)），则 \(h=\frac{1}{\sqrt2}f(b)\)，且 \(f'(b)=\frac{b}{\sqrt{b^2+1}}-1<0\)，故 \(h\) 最大于 \(b=0\)，无最小值，故面积有最大值而无最小值.
\end{solution}

\begin{problem}
设 \(\lambda\in\R\)，数列 \(\{a_n\}\)，\(\{b_n\}\)，\(\{c_n\}\) 的通项分别为 \(a_n=10n-9\)，\(b_n=2^n\)，\(c_n=\lambda a_n+(1-\lambda)b_n\). 若对任意 \(\lambda\in[0,1]\)，\(a_n\)，\(b_n\)，\(c_n\) 均能作为三角形的三边长，则满足条件的正整数 \(n\) 有（\quad）

\choices
    {1个}
    {3个}
    {4个}
    {无穷多个}
\end{problem}

\begin{answer}
B
\end{answer}

\begin{solution}
由
\(c_n=\lambda a_n+(1-\lambda)b_n\)，\(\lambda\in[0,1]\)，
可得
\(
\min\{a_n,b_n\}\le c_n\le \max\{a_n,b_n\}<a_n+b_n
\)
因此对任意 \(\lambda\in[0,1]\)，都有 \(c_n<a_n+b_n\).
要使 \(a_n\)，\(b_n\)，\(c_n\) 对任意 \(\lambda\in[0,1]\) 都能构成三角形，只需再保证较大的那一边恒小于另外两边之和.

令 \(f(x)=2^x-10x+9\)，则 \(f'(x)=2^x\ln2-10\). 因为 \(f'(x)\) 单调递增，且 \(f'(3)<0\)，\(f'(4)>0\)，
所以存在 \(x_0\in(3,4)\)，使得 \(f'(x_0)=0\). 于是 \(f(x)\) 先减后增，至多有两个零点.
又 \(f(1)>0\)，\(f(2)<0\)，\(f(5)<0\)，\(f(6)>0\)，所以 \(f(x)\) 恰有两个零点，分别位于 \((1,2)\) 与 \((5,6)\) 内. 于是对正整数 \(n\)，有
\(
10n-9\le 2^n \iff n=1 \text{\ 或\ } n\ge 6
\)
\(
10n-9\ge 2^n \iff n=2,3,4,5
\)

分两类讨论.

当 \(a_n\le b_n\) 时，即 \(n=1\) 或 \(n\ge 6\)，有 \(a_n\le c_n\le b_n\). 此时要构成三角形，只需 \(a_n+c_n>b_n\) 对任意 \(c_n\) 成立. 因为 \(c_n\) 的最小值为 \(a_n\)，所以只需
\(
2a_n>b_n
\)
即 \(2(10n-9)>2^n\). 与 \(a_n\le b_n\) 合并，只有 \(n=6\) 满足条件.

当 \(a_n\ge b_n\) 时，即 \(n=2\)，\(3\)，\(4\)，\(5\)，有 \(b_n\le c_n\le a_n\). 此时要构成三角形，只需 \(b_n+c_n>a_n\) 对任意 \(c_n\) 成立. 因为 \(c_n\) 的最小值为 \(b_n\)，所以只需
\(
2b_n>a_n
\)
即 \(2^{n+1}>10n-9\). 在 \(n=2\)，\(3\)，\(4\)，\(5\) 中，满足该不等式的只有 \(n=4\)，\(\ 5\).

综上，满足条件的正整数 \(n\) 为 \(4\)，\(5\)，\(6\)，共有 \(3\) 个，故选 B.
\end{solution}

\section{解答题}

\begin{problem}
2024 年东京奥运会，中国获得了男子 \(4\times100\) 米混合泳接力金牌，以下是历届奥运会男子 \(4\times100\) 米混合泳接力项目冠军成绩（单位：秒），数据按照升序排列.

\begin{center}
\begin{tabular}{|c|c|c|c|c|}
\hline
206.78 & 207.46 & 207.95 & 209.34 & 209.35 \\
\hline
210.68 & 213.73 & 214.84 & 216.93 & 216.93 \\
\hline
\end{tabular}
\end{center}

\begin{enumerate}
\item 求这组数据的极差与中位数；
\item 从这 10 个数据中任选 3 个，求恰有 2 个数据大于 \(211\) 的概率；
\item 若比赛成绩 \(y\) 关于年份 \(x\) 的回归方程为 \(y=-0.311x+\hat b\)，年份 \(x\) 的平均数为 \(2006\)，预测 \(2028\) 年男子 \(4\times100\) 米混合泳接力项目冠军队的成绩（精确到 \(0.01\) 秒）.
\end{enumerate}
\end{problem}

\begin{solution}
\begin{enumerate}
\item 数据的最小值为 \(206.78\)，最大值为 \(216.93\)，故极差 \(216.93-206.78=10.15\)；中间两数为 \(209.35,210.68\)，中位数为 \(\frac{209.35+210.68}{2}=210.015\).

\item 设恰有 2 个数据大于 \(211\)，事件 \(A\). 则从 \(4\) 个大于 \(211\) 的数中选两个，且从剩下 \(6\) 个中选一个：
\(P(A)=\frac{\mathrm C_4^2\mathrm C_6^1}{\mathrm C_{10}^3}=\frac{36}{120}=\frac3{10}\).

\item 先求样本均值
\[
\begin{aligned}
\bar y
&=\frac{206.78+207.46+207.95+209.34+209.35}{10}\\
&\quad+\frac{210.68+213.73+214.84+216.93+216.93}{10}=211.399.
\end{aligned}
\]
由 \((\bar x,\bar y)=(2006,211.399)\) 在回归线上，得 \(\hat b=835.265\)，故回归方程为
\(
y=-0.311x+835.265
\)
代入 \(x=2028\)，得
\(
y=-0.311\times2028+835.265 =204.557\approx204.56
\)
\end{enumerate}
\end{solution}

\begin{problem}
如图，\(P\) 是圆锥顶点，\(O\) 为底面圆心，\(AB\) 为底面直径，且 \(AB=2\).
\begin{enumerate}
\item 若直线 \(PA\) 与圆锥底面所成角为 \(\pi/3\)，求圆锥侧面积；
\item 已知 \(Q\) 是母线 \(PA\) 中点，点 \(C\)，\(D\) 在底面圆周上，且弧 \(AC\) 长为 \(\pi/3\)，\(CD\parallel AB\). 设点 \(T\) 在线段 \(OC\) 上，证明 \(QT\parallel PBD\) 所在平面.
\end{enumerate}

\bitmapfigure[max width=6.0cm]{img/2025/shanghai/q18_fig1.png}
\end{problem}

\begin{solution}
\begin{enumerate}
\item 由 \(\angle PAB=\pi/3\)，轴截面 \(\triangle ABP\) 为等边三角形，所以 \(PA=AB=2\). 底面半径 \(r=1\)，故母线 \(=2\). 侧面积 \(S=\pi r l=2\pi\).

\item 在三角形 \(PAB\) 中，\(Q\) 与 \(O\) 分别是 \(PA\) 与 \(AB\) 的中点，所以
\(QO\parallel PB\).
又 \(AB\perp AO\) 且 \(AB\parallel CD\). 设 \(\angle AOC=\pi/3\)，则 \(CD=1\)，且
\(OC\parallel BD\).
由于 \(QO\)，\(OC\) 是平面 \(\Pi_{QOC}\) 内的两条相交直线，\(PB\)，\(BD\) 是平面 \(\Pi_{PBD}\) 内的两条相交直线，且
\(QO\parallel PB\)，\(OC\parallel BD\)，
所以
\(\Pi_{QOC}\parallel\Pi_{PBD}\).
又 \(T\in OC\)，故 \(QT\subset\Pi_{QOC}\)，从而
\(QT\parallel\Pi_{PBD}\).
\end{enumerate}
\end{solution}

\begin{problem}
已知 \(f(x)=x^2-(m+2)x+m\ln x\)，\(m\in\R\).
\begin{enumerate}
\item 若 \(f(1)=0\)，求不等式 \(f(x)\le x^2-1\) 的解集；
\item 若 \(f(x)\) 在 \((0,+\infty)\) 上存在极大值，求 \(m\) 的取值范围.
\end{enumerate}
\end{problem}

\begin{solution}
\begin{enumerate}
\item 由 \(f(1)=0\) 得 \(m=-1\)，故
\(f(x)=x^2-x-\ln x\).
则
\(
f(x)\le x^2-1\iff x+\ln x\ge1
\)
设 \(s(x)=x+\ln x\)，\(x>0\)，\(s'(x)=1+\frac1x>0\)，故 \(s\) 在 \((0,+\infty)\) 上增，且 \(s(1)=1\)，解集为 \([1,+\infty)\).

\item 由
\[
f'(x)=2x-m-2+\frac{m}{x}=\frac{2x^2-(m+2)x+m}{x}=\frac{\left(2x-m\right)\left(x-1\right)}{x}
\]

当 \(m\le0\) 时，\(\frac{m}{2}\le0\)，在 \(x>0\) 上只有零点 \(x=1\). 当 \(x\in(0,1)\) 时 \(f'(x)<0\)，当 \(x\in(1,+\infty)\) 时 \(f'(x)>0\)，故 \(x=1\) 为极小值点，不符合题意.

当 \(0<m<2\) 时，\(0<\frac{m}{2}<1\). 当 \(x\in\left(0,\frac{m}{2}\right)\cup(1,+\infty)\) 时 \(f'(x)>0\)，当 \(x\in\left(\frac{m}{2},1\right)\) 时 \(f'(x)<0\)，故 \(x=\frac{m}{2}\) 为极大值点，符合题意.

当 \(m=2\) 时，\(f'(x)=\frac{2\left(x-1\right)^2}{x}\ge0\)，故 \(f(x)\) 无极大值点，不符合题意.

当 \(m>2\) 时，\(\frac{m}{2}>1\). 当 \(x\in(0,1)\cup\left(\frac{m}{2},+\infty\right)\) 时 \(f'(x)>0\)，当 \(x\in\left(1,\frac{m}{2}\right)\) 时 \(f'(x)<0\)，故 \(x=1\) 为极大值点，符合题意.

综上，\(m\) 的取值范围为 \(m>0\) 且 \(m\ne2\).
\end{enumerate}
\end{solution}

\begin{problem}
已知椭圆 \(\Gamma:\frac{x^2}{a^2}+\frac{y^2}{5}=1\)（\(a>\sqrt5\)），\(M(0,m)\)（\(m>0\)），\(A\) 为 \(\Gamma\) 的右顶点.
\begin{enumerate}
\item 若 \(\Gamma\) 的焦点为 \((2,0)\)，求离心率 \(e\)；
\item 若 \(a=4\)，且 \(\Gamma\) 上存在一点 \(P\) 满足 \(\overrightarrow{PA}=2\overrightarrow{MP}\)，求 \(m\)；
\item 若线段 \(AM\) 的中垂线 \(l\) 的斜率为2， \(l\) 与 \(\Gamma\) 交于 \(C\)，\(D\)，且 \(\angle CMD\) 为钝角，求 \(a\) 的取值范围.
\end{enumerate}
\end{problem}

\begin{solution}
\begin{enumerate}
\item 由 \(b^2=5\)，焦点 \((\pm c,0)\) 有 \(c=2\)，所以 \(a=\sqrt{c^2+b^2}=\sqrt9=3\)，故 \(e=\frac{c}{a}=\frac23\).

\item 此时 \(A(4,0)\)，设 \(P(x_P,y_P)\). 由 \(\overrightarrow{PA}=2\overrightarrow{MP}\) 得
\(4-x_P=2x_P\)，\(-y_P=2(y_P-m)\).
所以 \(x_P=\frac43\)，\(y_P=\frac{2m}{3}\). 代入椭圆方程 \(\frac{x^2}{16}+\frac{y^2}{5}=1\) 得 \(\frac{1}{9}+\frac{4m^2}{45}=1\)，故 \(m=\sqrt{10}\).

\item 由 \(l\) 为 \(AM\) 的中垂线且斜率为 \(2\)，知 \(AM\) 的斜率为 \(-\frac12\). 又 \(k_{AM}=\frac{m-0}{0-a}=-\frac{m}{a}\)，故 \(m=\frac a2\). 于是 \(AM\) 的中点为 \(\left(\frac a2,\frac a4\right)\)，直线 \(l\) 的方程为 \(y-\frac a4=2\left(x-\frac a2\right)\)，即 \(y=2x-\frac{3a}{4}\).
联立
\(
\begin{cases}
\frac{x^2}{a^2}+\frac{y^2}{5}=1,\\
y=2x-\frac{3a}{4},
\end{cases}
\)
得
\((4a^2+5)x^2-3a^3x+\frac{9a^4-80a^2}{16}=0\).
设两交点横坐标为 \(x_1\)，\(x_2\)，则由韦达定理可得
\(x_1+x_2=\frac{3a^3}{4a^2+5}\)，\(x_1x_2=\frac{a^2(9a^2-80)}{16(4a^2+5)}\).
又 \(M\left(0,\frac a2\right)\)，且 \(y_i=2x_i-\frac{3a}{4}\)（\(i=1\)，\(2\)）. \(\angle CMD\) 为钝角当且仅当 \(\overrightarrow{MC}\cdot\overrightarrow{MD}<0\)，即 \(x_1x_2+\left(y_1-\frac a2\right)\left(y_2-\frac a2\right)<0\).
代入 \(y_i=2x_i-\frac{3a}{4}\) 得
\(
5x_1x_2-\frac{5a}{2}(x_1+x_2)+\frac{25a^2}{16}<0
\)
再代入韦达定理，化简得
\(
\frac{25a^2(a^2-11)}{16(4a^2+5)}<0
\)
因为 \(a>\sqrt5\)，所以 \(16(4a^2+5)>0\)，从而 \(a^2<11\). 再结合 \(a>\sqrt5\)，得
\(
\sqrt5<a<\sqrt{11}
\)
\end{enumerate}
\end{solution}

\begin{problem}
已知函数 \(y=f(x)\) 定义域为 \(\R\). 对于正实数 \(a\)，定义集合 \(M_a=\{x\mid f(x+a)=f(x)\}\).
\begin{enumerate}
\item 若 \(f(x)=\sin x\)，判断 \(\dfrac{\pi}{3}\) 是否是 \(M_{\pi}\) 的元素，请说明理由；
\item 若 \(f(x)=x+2\)（\(x<0\)），\(f(0)=0\)，且当 \(x>0\) 时 \(f(x)=\sqrt{x}\)，\(M_a\neq\varnothing\)，求 \(a\) 的取值范围；
\item 若 \(y=f(x)\) 是偶函数，当 \(x\in(0,1]\) 时 \(f(x)=1-x\)，且对任意 \(a\in(0,2)\)，有 \(M_a\subset M_2\). 写出 \(f(x)\) 在 \(x\in(1,2)\) 的解析式，并证明：对任意实数 \(c\)，函数 \(f(x)-c\) 在 \([-3,3]\) 上至多有9个零点.
\end{enumerate}
\end{problem}

\begin{solution}
\begin{enumerate}
\item 计算得 \(f\left(\frac{\pi}{3}\right)=\sin\frac{\pi}{3}=\frac{\sqrt{3}}{2}\)，且 \(f\left(\frac{\pi}{3}+\pi\right)=-\sin\frac{\pi}{3}=-\frac{\sqrt{3}}{2}\). 所以 \(\dfrac{\pi}{3}\notin M_{\pi}\).

\item 取 \(x_0\in M_a\)，则 \(f(x_0+a)=f(x_0)\) 且 \(a>0\).

若 \(x_0<0\) 且 \(x_0+a<0\)，则 \(f(x_0+a)=x_0+a+2\ne x_0+2=f(x_0)\)，矛盾；若 \(x_0\ge0\)，则 \(x_0+a>0\)，从而 \(f(x_0+a)=\sqrt{x_0+a}\ne\sqrt{x_0}=f(x_0)\)，也矛盾.

故必有 \(x_0<0\le x_0+a\). 于是 \(x_0+2=\sqrt{x_0+a}\).

因为 \(-2\le x_0<0\)，所以
\(
a=(x_0+2)^2-x_0=\left(x_0+\frac{3}{2}\right)^2+\frac{7}{4}\in\left[\frac{7}{4},4\right)
\)
故 \(a\in\left[\frac{7}{4},4\right)\).

\item 任取 \(x_0\in(1,2)\)，则 \(x_0-2\in(-1,0)\). 由于 \(f\) 为偶函数，\(f(x_0-2)=f(2-x_0)\). 又 \((2-x_0)-(x_0-2)=4-2x_0\in(0,2)\)，所以 \(x_0-2\in M_{4-2x_0}\subset M_2\)，从而
\[
f(x_0)=f(x_0-2)=f(2-x_0)=1-(2-x_0)=x_0-1
\]
故当 \(x\in(1,2)\) 时，\(f(x)=x-1\).
当 \(s\in(0,1)\) 时，\(1-s\in(0,1)\)，\(1+s\in(1,2)\)，于是 \(f(1-s)=1-(1-s)=s\)，\(f(1+s)=(1+s)-1=s\). 故 \(1-s\in M_{2s}\subset M_2\)，从而 \(f(3-s)=f(1-s)=s\). 由于 \(3-s\in(2,3)\)，所以当 \(x\in(2,3)\) 时，\(f(x)=3-x\).

再由偶性可得
\(f(x)=x+1\)，\(x\in(-1,0)\)，\(f(x)=-x-1\)，\(x\in(-2,-1)\)，\(f(x)=x+3\)，\(x\in(-3,-2)\).

设 \(f(-3)=n\). 若 \(n\in(0,1)\)，则 \(-3+n\in(-3,-2)\)，从而 \(f(-3+n)=(-3+n)+3=n=f(-3)\)，故 \(-3\in M_n\subset M_2\)，于是 \(f(-1)=f(-3)=n\).
但由偶性及 \(f(1)=0\)，得 \(f(-1)=0\)，矛盾. 所以 \(f(-3)\notin(0,1)\). 再由偶性，得 \(f(3)\notin(0,1)\).

令 \(f(x)-c=0\)，即 \(f(x)=c\).

若 \(c\notin(0,1)\)，则六个开区间
\((-3,-2)\)，\((-2,-1)\)，\((-1,0)\)，\((0,1)\)，\((1,2)\)，\((2,3)\)
内均无零点，此时零点只可能出现在 \(-3\)，\(-2\)，\(0\)，\(2\)，\(3\) 这 5 个点上；特别地，当 \(c=0\) 时，又有 \(x=\pm1\) 两个零点，所以此时零点至多有 7 个.

若 \(0<c<1\)，则上述 6 个开区间内各至多有 1 个零点；又 \(f(\pm3)\notin(0,1)\)，所以只可能在 \(-2\)，\(0\)，\(2\) 这 3 个点处再出现零点，从而零点总数不超过 \(6+3=9\).

综上，函数 \(f(x)-c\) 在 \([-3,3]\) 上至多有 9 个零点.
\end{enumerate}
\end{solution}
